\documentclass[11pt,a4paper]{article}
\usepackage[utf8]{inputenc}
\usepackage[T1]{fontenc}
\usepackage[brazil]{babel}
\usepackage[a4paper,margin=2.3cm]{geometry}
\usepackage{enumitem}
\usepackage{hyperref}
\usepackage{xcolor}
\setlength{\parindent}{0pt}
\setlength{\parskip}{5pt}
\sloppy
\setlength{\emergencystretch}{3em}
\hypersetup{colorlinks=true,urlcolor=blue,linkcolor=black}
\newcommand{\cabecalho}[2]{\begin{center}{\small \textbf{Sprint 2 -- #1}}\\[1mm]\textit{Inventário de Ativos e Vulnerabilidades}\\Matheus Moreli Durão Silva -- 12621CBS217\\Cibersegurança -- Bacharelado -- Noturno\end{center}\hrule\vspace{4mm}}
\newcommand{\bloco}[1]{\par\vspace{2mm}\textbf{#1}\par}
\newcommand{\codigo}[1]{\texttt{#1}}

\begin{document}\cabecalho{S2\_09 -- Busca de Ativos}{}
\bloco{Tipos de busca} A busca por ID usa \codigo{consultar\_por\_id} e devolve um único ativo ou \codigo{None}. A busca com filtros usa \codigo{filtrar\_ativos}: termo parcial do nome, tipo e criticidade podem ser combinados.
\bloco{Como a filtragem funciona} O \codigo{for} percorre os objetos do dicionário. Para cada ativo que não atende a um filtro, \codigo{continue} pula para o próximo. Quem atende é acrescentado a uma nova lista chamada \codigo{resultado}. A base original nunca recebe modificação.
\bloco{Exemplo} Buscar termo \codigo{srv} e criticidade \codigo{Alta} encontra \codigo{SRV-ARQUIVOS} se os dois critérios forem verdadeiros. Uma busca sem resultado mostra mensagem previsível; não tenta acessar uma chave inexistente.
\bloco{Como testar} O teste \codigo{testar\_filtro\_nao\_altera\_a\_base} confirma simultaneamente que há um resultado e que a quantidade da base permanece igual.
\bloco{Checklist atendido} ID único, filtros previsíveis, consulta sem alteração, ausência tratada e casos cobertos por teste.
\end{document}
